%%%%%%%%%%%%%%%%%%%%%%%%%%%%%%%%%%%%%%%%%%%%%%%%%%%%%%%%%%
% FILE: chapter08_build_mode.tex
%
% PURPOSE
% -------------------------------------------------------
% Explains the build modes provided by the template.
%
%
% ARCHITECTURE ROLE
% -------------------------------------------------------
% User guide chapter.
%
% Documents the engineering build system used to
% control draft and final compilation modes.
%
%
% DEPENDENCIES
% -------------------------------------------------------
% config.tex
% packages.tex
% tikz
%
%
% NOTES FOR DEVELOPERS
% -------------------------------------------------------
% If the build system changes, this chapter must be
% updated accordingly.
%
%%%%%%%%%%%%%%%%%%%%%%%%%%%%%%%%%%%%%%%%%%%%%%%%%%%%%%%%%%


\chapter{Modo build}

La plantilla incorpora un sistema de compilación que
permite generar diferentes versiones del documento.

Este sistema facilita el desarrollo del documento y
la generación de la versión final.



% --------------------------------------------------
% CONCEPTO DE BUILD
% --------------------------------------------------

\section{Concepto de build}

\begin{definitionbox}[Build]
Un \textit{build} es el proceso mediante el cual el
código fuente de un proyecto se transforma en su
resultado final.
\end{definitionbox}

En el caso de \LaTeX{}, el build corresponde al
proceso de compilación que genera el archivo PDF.



% --------------------------------------------------
% MODOS DE COMPILACIÓN
% --------------------------------------------------

\section{Modos de compilación}

La plantilla define dos modos principales:

\begin{itemize}

\item modo \textbf{draft}
\item modo \textbf{final}

\end{itemize}



% --------------------------------------------------
% MODO DRAFT
% --------------------------------------------------

\section{Modo draft}

El modo \texttt{draft} está diseñado para el proceso
de escritura y revisión del documento.

\begin{lstlisting}[caption={Activar modo draft},label={code:draftmode}]
\drafttrue
\end{lstlisting}

\source{Plantilla}



En este modo el documento puede mostrar información
adicional útil durante el desarrollo.



\begin{examplebox}
Ejemplos de información útil en modo draft:

\begin{itemize}

\item información de compilación
\item fecha de generación del documento
\item identificadores internos del documento

\end{itemize}

\end{examplebox}



% --------------------------------------------------
% MODO FINAL
% --------------------------------------------------

\section{Modo final}

El modo final está pensado para la versión definitiva
del documento.

\begin{lstlisting}[caption={Modo final},label={code:finalmode}]
\draftfalse
\end{lstlisting}

\source{Plantilla}



En este modo el documento se genera sin información
de depuración ni elementos auxiliares.



\begin{notebox}
Por defecto la plantilla compila en modo final.
\end{notebox}



% --------------------------------------------------
% INFORMACIÓN DE BUILD
% --------------------------------------------------

\section{Información de compilación}

La plantilla puede mostrar información sobre el
proceso de compilación.

\begin{lstlisting}[caption={Macro de información de build},label={code:buildinfo}]
\BuildInfo
\end{lstlisting}

\source{Plantilla}



Esta información puede incluir:

\begin{itemize}

\item fecha de compilación
\item versión del documento
\item estado del build

\end{itemize}



% --------------------------------------------------
% METADATOS DEL DOCUMENTO
% --------------------------------------------------

\section{Metadatos del documento}

El archivo PDF generado incluye metadatos que
describen el documento.

Estos metadatos se configuran mediante:

\begin{lstlisting}[caption={Metadatos del PDF},label={code:metadata}]
\hypersetup{
  pdftitle={\documenttitle},
  pdfauthor={\documentauthor},
  pdfcreator={mammesa}
}
\end{lstlisting}

\source{Plantilla}



\begin{definitionbox}[Metadatos]
Los metadatos son información interna del archivo
PDF que describe el contenido del documento.
\end{definitionbox}



% --------------------------------------------------
% LOG DE COMPILACIÓN
% --------------------------------------------------

\section{Log de compilación}

Durante el proceso de compilación, la plantilla
puede generar mensajes informativos en el log.

\begin{lstlisting}[caption={Mensaje de compilación},label={code:logmessage}]
\typeout{mammesa-template-loaded}
\end{lstlisting}

\source{Plantilla}



Este tipo de mensajes permite verificar que los
componentes principales de la plantilla se han
cargado correctamente.



\begin{warningbox}
Los mensajes del log no aparecen en el documento
final, pero pueden ser útiles para depurar
problemas de compilación.
\end{warningbox}



% --------------------------------------------------
% VENTAJAS DEL SISTEMA
% --------------------------------------------------

\section{Ventajas del sistema}

El sistema de build ofrece varias ventajas:

\begin{itemize}

\item separación entre documento en desarrollo y final
\item control centralizado de la compilación
\item información adicional para depuración
\item metadatos estructurados del documento

\end{itemize}



\begin{examplebox}
Este enfoque está inspirado en sistemas de
construcción utilizados en proyectos de software,
adaptados al flujo de trabajo de \LaTeX.
\end{examplebox}
